\par\Needspace{6\baselineskip}\medskip
\SourceLeaf{195}\begin{center}\hypertarget{native-195-article144}{}\pdfbookmark[2]{144.}{article-144}144.\end{center}
Once the arcs $\zeta,\zeta''$ have been found, the positions of $C,C''$ are given, and the distance $CC''=2f'$ can be determined from $\zeta,\zeta''$ and $\epsilon'$. Let $u,u''$ be the inclinations of the great circles $AB,A''B''$ to the great circle $CC''$ (in Fig. 4 these are the angles $C''CD'$ and $180^{0}-CC''D'$ respectively). We then have the following equations, exactly analogous to equations 3—6 of Article 137.
\[
\begin{aligned}
\sin f'\sin\tfrac12(u''+u)&=\sin\tfrac12\epsilon'\sin\tfrac12(\zeta+\zeta'')\\
\sin f'\cos\tfrac12(u''+u)&=\cos\tfrac12\epsilon'\sin\tfrac12(\zeta-\zeta'')\\
\cos f'\sin\tfrac12(u''-u)&=\sin\tfrac12\epsilon'\cos\tfrac12(\zeta+\zeta'')\\
\cos f'\cos\tfrac12(u''-u)&=\cos\tfrac12\epsilon'\cos\tfrac12(\zeta-\zeta'').
\end{aligned}
\]
The first two give $\tfrac12(u''+u)$ and $\sin f'$, and the last two give $\tfrac12(u''-u)$ and $\cos f'$. From $\sin f'$ and $\cos f'$ we obtain $f'$. The angles $\tfrac12(u''+u)$ and $\tfrac12(u''-u)$ are used only in the final hypothesis to determine the position of the plane of the orbit, and may be omitted in the initial hypotheses.

In exactly the same way, $f$ could be derived from $\epsilon,C'D$ and $C''D$, and $f''$ from $\epsilon'',CD'',C'D''$. But the following formulas are much more convenient for this purpose.
\[
\begin{aligned}
\sin2f&=r\sin2f'\cdot\frac{n}{n'r'}\\
\sin2f''&=r''\sin2f'\cdot\frac{n''}{n'r'}\mathpunct{\raisebox{.45ex}{,}}
\end{aligned}
\]
Here the logarithms of $\dfrac{n}{n'r'}\mathpunct{\raisebox{.45ex}{,}}$ $\dfrac{n''}{n'r'}$ are already available from the preceding calculations. Finally, the entire calculation has a further check, since we must have $2f+2f''=2f'$. Any difference that does appear can certainly be of no consequence if all the operations have been carried out as accurately as possible. Sometimes, however, even when the calculation is carried out throughout to seven decimal places, it can amount to a few tenths of an arcsecond. If it seems worthwhile, we can very easily distribute the difference between $2f,2f''$ so that the logarithms of the sines are both increased or both decreased by the same amount. This satisfies the equation $P=\dfrac{r\sin2f''}{r''\sin2f}=\dfrac{n''}{n}$ to the full precision allowed by the tables. When $f$ and $f''$ differ little, it suffices to distribute the difference equally between $2f$ and $2f''$.
